% !TEX root = ../main.tex
% Draft 10: evaluation skeleton aligned to the revised claim; no new experiment has been run.
\section{Experiments}
\label{sec:experiments}
The evaluation is organized around five questions:
\textbf{Q1}, how task completion and forbidden contact change;
\textbf{Q2}, what attention contributes to target identity and geometric extent;
\textbf{Q3}, what phase information adds beyond grasp-state rules;
\textbf{Q4}, which contact representation matters under a fixed controller; and
\textbf{Q5}, what the complete system costs at run time.
\tbd{Finalize the evaluation protocol and replace this skeleton with registered results.}

\subsection{Experimental setup}
\paragraph{Tasks and policy.}
The existing implementation uses a $\pi_{0.5}$ checkpoint fine-tuned for dual-arm pick-and-place in MuJoCo~\cite{todorov2012mujoco},
with one supplied active manipulator per episode. The diagnostic scenes place distractors near the grasp target and
obstacles along the carrying path. Simulation uses rendered depth and waits for the inference server.
\tbd{Fix the evaluated snapshot, scene conditions, episodes, seeds, and training/tuning/test separation.}
SafeLIBERO~\cite{hu2026vlsa} is the planned primary benchmark, while the RB-Y1 simulation supports baseline comparisons
and component ablations. A real RB-Y1 demonstration is planned to assess execution feasibility, without a comparative
performance claim. \tbd{Complete the SafeLIBERO observation/action adapter, fix its checkpoint and evaluation suites,
and document whether phase handling uses HSR or a grasp-state rule. Keep robot demonstrations marked pending until recorded.}

\paragraph{Comparisons.}
On SafeLIBERO, we plan direct runs of $\pi_{0.5}$ alone and with \sys, and comparisons with two or three of VLSA,
KNOWS, and WBAG~\cite{hu2026vlsa,park2026knows,zhen2026wbag}. Published values retain their source metric definitions;
incompatible checkpoints, scenarios, or metrics require separate reporting blocks. On RB-Y1, available baseline code
will be adapted, or a reimplementation explicitly identified, prioritizing KNOWS and VLSA. Component comparisons
instead retain our controller and replace one target source, phase rule, or contact representation at a time.
\tbd{Confirm baseline versions, reported versus rerun provenance, input parity, and tuning budgets.}

\paragraph{Metrics.}
We will report task success, the fraction of trials containing a forbidden contact, and their joint metric,
\emph{safe success}. Allowed contacts must be specified by object pair, robot link, and task stage; the target and
destination are not blanket exclusions. Contact scoring uses simulator records independently of the ESDF being tested.
\tbd{Freeze success/release criteria, the allowed-contact table, sampled-time coverage, and treatment of stops and
incomplete trials before evaluation.} Safe success is computed per trial, not by multiplying aggregate rates.
We also report target identity accuracy, extent quality and obstacle retention, stop frequency, and end-to-end latency.
\tbd{Specify sample counts, aggregation units, and paired confidence intervals.}

\begin{table*}[t]
\centering\small
\begin{tabularx}{\textwidth}{@{}>{\raggedright\arraybackslash}Xllccc@{}}
\toprule
Method & Result source & Query coverage & Success $\uparrow$ & Contact $\downarrow$ & Safe $\uparrow$ \\
\midrule
$\pi_{0.5}$ only & planned rerun & no added filter & -- & -- & -- \\
\sys\ (ours) & planned rerun & \tbd{final links} & -- & -- & -- \\
\midrule
VLSA~\cite{hu2026vlsa} & \tbd{reported/rerun} & end effector & -- & -- & -- \\
KNOWS~\cite{park2026knows} & \tbd{reported/rerun} & end effector & -- & -- & -- \\
WBAG~\cite{zhen2026wbag} & \tbd{reported/rerun} & whole body + attached & -- & -- & -- \\
\bottomrule
\end{tabularx}
\caption{SafeLIBERO main-comparison skeleton. The lower block lists baseline candidates, not completed runs.
Values can share metric columns only after definitions and conditions are matched; incompatible published metrics
must receive separate labeled blocks. Privileged simulator input in WBAG must be identified in the final comparison.
\tbd{Select baselines; insert suite/level, sample counts, exact metrics, result sources, and uncertainty.}}
\label{tab:main}
\end{table*}

\begin{table}[t]
\centering\small
\begin{tabularx}{\columnwidth}{@{}>{\raggedright\arraybackslash}Xccc@{}}
\toprule
Change from full system & Success $\uparrow$ & Contact $\downarrow$ & Safe $\uparrow$ \\
\midrule
Geometry-only target selection & -- & -- & -- \\
External-segmenter target & -- & -- & -- \\
No phase gate & -- & -- & -- \\
No target exemption & -- & -- & -- \\
Global rather than link-specific exclusion & -- & -- & -- \\
No attached-object queries & -- & -- & -- \\
Grasp-state rule rather than HSR & -- & -- & -- \\
Full system & -- & -- & -- \\
\bottomrule
\end{tabularx}
\caption{Matched-controller ablation skeleton. Each row changes one factor with other settings fixed.
\tbd{Finalize source/gate definitions, hold attachment settings fixed in the HSR comparison, and insert measured results.}}
\label{tab:ablation}
\end{table}

\subsection{Results and analysis}
\paragraph{Q1: task completion and forbidden contact.}
\tbd{Use Table~\ref{tab:main} to state the supported task/contact trade-off. Explain baseline failures using recorded
contact events and matched rollouts; do not infer physical safety from fewer stops. Add a separate RB-Y1 simulation
comparison table once baseline adapters are fixed, and a real-robot demonstration panel if execution is completed.}

\paragraph{Q2: target identity and extent.}
\tbd{Compare ATS with geometry-only selection and an external segmenter under the same contact rules and CCR.
Measure target identity separately from observed extent and wrongly excluded obstacle geometry.
Figure~\ref{fig:fig03_attention_to_3d} should connect these measurements to changed constraints and execution.}

\paragraph{Q3: phase-dependent target handling.}
\tbd{Compare HSR with a grasp-state/history rule using the same ATS, attachment logic, and controller.
Report downstream errors such as stale exceptions and erroneous target switching, with task/contact outcomes.
Use independently defined phase labels; determine whether the benefit supports a main-text or supplementary claim.}

\paragraph{Q4: contact representation.}
\tbd{Use Table~\ref{tab:ablation} to separate global exclusion, link-specific queries, and attached geometry.
Pair each supported effect with a recorded failure mechanism; retain only comparisons needed for the central claim.}

\paragraph{Q5: computational cost.}
\tbd{Report policy serving, separate attention extraction, probe, ATS, field construction, CCR, and communication time.
Measure end-to-end latency against the nominal 533 ms execution horizon; provide the detailed breakdown in Appendix D.}
